\section{相干态泛函积分}

对多体体系，用相干态闭合关系代替坐标完备性，得到场的泛函积分。

\subsection{玻色子}

虚时作用量（正常序哈密顿量）
\begin{equation}
  S[\phi^*,\phi]
  =\int_0^{\hbar\beta}\mathrm{d}\tau
  \Bigl[
  \sum_\alpha\phi_\alpha^*\hbar\partial_\tau\phi_\alpha
  +H(\phi^*,\phi)
  \Bigr],
\end{equation}
配分函数
\begin{equation}
  Z
  =\int\mathcal{D}[\phi^*,\phi]\,
  e^{-S/\hbar},
\end{equation}
场满足周期边界 $\phi(\hbar\beta)=\phi(0)$。

\subsection{费米子}

用 Grassmann 场 $\psi,\psi^*$，作用量形式类似，但
\begin{equation}
  \psi(\hbar\beta)=-\psi(0)
\end{equation}
（反周期），与 Matsubara 费米频率 $(2n+1)\pi/\beta$ 一致。二次型作用量的高斯积分直接给出 $\det$，对应自由费米子配分函数。

\subsection{与二次量子化的等价性}

泛函积分不引入新物理：它把 $\mathrm{Tr}\,T_\tau\exp(-\int H)$ 改写为场积分。优点是：
\begin{itemize}
  \item 平均场 = 鞍点；
  \item 涨落 = 高斯积分 / 环图；
  \item Hubbard--Stratonovich 变换可引入辅助场（密度、配对、自旋等）。
\end{itemize}
